\documentclass[11pt,a4paper,toc=bibliography,toc=listof,titlepage=firstiscover]{scrreprt}
\usepackage[utf8]{inputenc}
\usepackage{blindtext}
\usepackage{graphicx}
\usepackage[T1]{fontenc}
\usepackage{setspace}
\usepackage{pdfsync}
\usepackage{blindtext}
\usepackage[ngerman]{babel}
\usepackage{booktabs}
\usepackage{xcolor}
\usepackage{multirow}
%Schriftarten hier definieren
\usepackage{lmodern}
%\usepackage{tgheros}
%\renewcommand*\familydefault{\sfdefault}
%\usepackage{tgpagella}
\usepackage[babel,german=guillemets]{csquotes}
\usepackage[style=ieee-alphabetic]{biblatex}
\usepackage{hyperref}
\usepackage{csquotes}
\addbibresource{literatur.bib}
\hypersetup{colorlinks=true, linkcolor=black, urlcolor=black, citecolor=black}
%Abstand der Einträge im Literaturverzeichnis als em
\setlength\bibitemsep{0.75em}
%Kein Einzug bei einem neuen Absatz
\setlength{\parindent}{0mm}
%Inhaltsverzeichnis ohne eigene Seitenzahl
\AtBeginDocument{\addtocontents{toc}{\protect\thispagestyle{empty}}} 
\usepackage[left=3cm,right=2.5cm,top=2.5cm,bottom=2cm]{geometry}

%Daten des Deckblatts
  \titlehead
  {
    {
      \begin{center}
        \vspace{2.5cm}
        \includegraphics[scale=0.3]{resources/nak-logo.png}
      \end{center}
    }
  	%M.Sc. Angewandte Informatik
  }
  \subject{Hausarbeit}
  \title{Integration der Data Science in Organisationen / Abteilungen}
    \subtitle{Vorlesung Data-Science Projektorganisation}
    \author{Leon Henne}
  \date{\small{Köln, den \today}}
  \publishers{Betreut durch Prof. Dr. Michael Schulz}
\begin{document}
%Titelseite erzeugen
\maketitle
%Der Nachfolgende Text wird 1,5-zeilig gesetzt
\onehalfspacing
%Römische Ziffern für Abbildungsverzeichnis
\pagenumbering{Roman}
\setcounter{page}{0}
%Inhaltsverzeichnis erzeugen
\tableofcontents
\clearpage
%Abbildungs und Tabellenverzeichnis
\listoftables
\listoffigures
%ToDo --> Abkürzungsverzeichnis
\clearpage
\pagenumbering{arabic}
\setcounter{page}{1}
%Datei aus dem Ordner /chapter/ einbinden. So kann man mehrere Textdateien zu einem großen Dokument zusammenfügen.
\setcounter{chapter}{0}
\chapter[Problemstellung]{Problemstellung}

Viele Trends der technologischen Welt innerhalb der letzten Jahrzehnte wurden besonders durch drei bestimmende Faktoren beeinflusst.
Die Entwicklung der Datenspeicherungs- und Rechenkapazitäten als zwei dieser Faktoren entstammen den Erkenntnissen und sich bestätigenden Entwicklungen von Gorden E. Moore bzw. dem Moorschen Gesetz. \footcite[Vgl.][S. 1]{Moore.1998}
Demnach wurde bereits sehr früh erkannt, dass integrierte Schaltungen die bekannten Telefonschaltungen ersetzen wird und Computer leistungsfähiger entsprechende Daten verarbeiten können. \footcite[Vgl.][S. 1]{Moore.1998}
Erkenntlich wird diese Entwicklung anhand der heutzutage konstanten und umfangreichen Generierung von Daten durch Anwendungen in Mobiltelefonen, Autos, IOT Geräten oder Industriemaschinen. \footcite[Vgl.][S 3f.]{Dalpiaz.2020}
Prognosen sagen dabei eine verfünffachende Entwicklung der globalen Datenmenge vorher von 33 Zettabytes in 2018 zu 175 Zettabytes in 2025. \footcite[Vgl.][S. 1]{Hupperz.2021}
Der dritte Faktor ist die Entwicklung der Möglichkeiten zur Datenverarbeitung und Analyse anhand von zunehmend komplexeren Algorithmen. \footcite[Vgl.][S. 4]{Dalpiaz.2020}

Folgend aus diesen drei Faktoren zielen immer mehr Software Unternehmen darauf ab, sich in eine datengesteuerte Organisation zu transformieren. \footcite[Vgl.][S. 1]{Fabijan.2017}
Höheres Interesse an dem Besitz und der Verarbeitung großer Datenmengen zeigen heutzutage jedoch Organisationen aller Bereiche wie Wirtschaft, Regierungen und Forschung. \footcite[Vgl.][S. 1]{Pratt.2023}
Deren Stratgien weisen häufig auf, mittels datenfokussierter Kultur und Datenanalysen möglichst umfangreiche faktenbasierte Entscheidungen zu treffen. \footcite[Vgl.][S. 18]{Dalpiaz.2020}
Forschungsfeld dieser Absicht bildet die Data Science, bzw. Big Data Analytics mit dessen Einsatz von Technologien sich Unternehmen einen Wettbewerbsvorteil erzielen. \footcite[Vgl.][S. 3]{Dalpiaz.2020}
datengesteuerte Unternehmen können dadurch bis zu 26 \% profitabler werden gegenüber Wettbewerbern, welche weniger bis gar keine digitalen Technologien einsetzen. \footcite[Vgl.][S. 1]{Fabijan.2017}
Ein Grund dafür können bessere und schnellere Entscheidungen sein, welche aus dem vollumfänglicheren Verständnis des Kunden und höherer Transparenz des Entwicklungsprozesses resultieren. \footcite[Vgl.][S. 18]{Dalpiaz.2020}
Trotz der heute vorhandenen Fülle von erzeugten Daten bleibt die Anzahl an Unternehmen, welche erfolgreich in eine datengesteuerte Organisation transformierten eher gering. \footcite[Vgl.][S. 1]{Fabijan.2017}
Zusätzlich existiert bisher nur wenig Forschung dazu, wie solche Transformationen zu datengesteuerte Unternehmen umzusetzen sind. \footcite[Vgl.][S. 1]{Fabijan.2017}

% \begin{itemize}
%     \item Entwicklung von Datenspeicherung, Rechenpower und Analysemöglichkeiten
%     \item Relevanz der Thematik für Organisationen als Wettbewerbsvorteil (was bringt es ?) 
%     \begin{itemize}
%         \item Steigende Notwendigkeit der Integration von Data Science zur Marktpositionierung
%     \end{itemize}
%     \item Aktualität des Problems (Bisher hat es halt keiner gemacht)
% \end{itemize}
\chapter[Grundlagen]{Grundlagen}

\section{Disziplin der Data Science}



\section{Organisationen}

\section{Abgrenzung von Daten getriebenen Organisationen}
\chapter[Integrationsprozesse]{Integrationsprozesse}

\section{Herausforderungen}

\section{Rollen in datengetriebenen Organisationen}

\section{Vorgehensmodelle}

\begin{itemize}
    \item Experiment Evolution Model (Microsoft)
    \item ???
\end{itemize}
\chapter[Betrachtung von Fallstudien]{Betrachtung von Fallstudien}
\chapter[Fazit]{Fazit}
%Literaturverzeichnis einfügen
\renewcommand*{\UrlFont}{\rmfamily}
\printbibliography
\include{chapter/99_eidesstattliche-erkläerung.tex}
\end{document}
